% ============================================================
% 《算法的低语·数学卷》独立分册 preamble
% 基于主册 preamble.tex 精简（去除主册封面与变更日志，新增数学卷封面）
% ============================================================

\usepackage{fontspec}
\usepackage[fontset=fandol]{ctex}
\usepackage{geometry}
\geometry{
  a4paper,
  top=2.7cm, bottom=2.7cm, left=2.8cm, right=2.8cm,
  headheight=1.2cm, footskip=1.4cm
}
\usepackage{amsmath, amssymb, mathtools}
\usepackage{graphicx}
\usepackage{float}
\usepackage{xcolor}
\usepackage{hyperref}
\usepackage{eso-pic}
\usepackage{tikz}
\usepackage{fancyhdr}
\usepackage{titlesec}
\usepackage{enumitem}
\usepackage{booktabs}
\usepackage{array}
\usepackage{calc}
\usepackage[most]{tcolorbox}

% ---- 颜色（与主册一致）----
\definecolor{teal1}{RGB}{0,77,89}
\definecolor{teal2}{RGB}{15,118,135}
\definecolor{teal3}{RGB}{38,162,180}

% ---- 占位符：用三角形作为导读盒标题前缀 ----
\newcommand{\faIconStub}{\textbf{\textcolor{teal2}{▶}}}

% ---- 小节导读盒：「这一节在讲什么」----
\newtcolorbox{sectionintuition}{
  enhanced,
  colback=teal3!6,
  colframe=teal2,
  boxrule=0.6pt,
  arc=3pt,
  left=10pt, right=10pt, top=8pt, bottom=8pt,
  before skip=8pt, after skip=10pt,
  fonttitle=\bfseries\color{teal1},
  title={\faIconStub\,这一节在讲什么},
  attach title to upper,
  after title={\par\vspace{4pt}}
}

% ---- 超链接 ----
\hypersetup{
  colorlinks=true,
  linkcolor=teal2,
  urlcolor=teal2,
  citecolor=teal2,
  pdftitle={算法的低语·数学卷习题伴读 v1.2 全册},
  pdfauthor={大队长}
}

% ---- 页眉页脚（与主册风格统一）----
\pagestyle{fancy}
\fancyhf{}
\fancyhead[L]{\small\color{gray}\itshape 算法的低语·数学卷}
\fancyhead[R]{\small\color{gray}\itshape 从公式到推导}
\fancyfoot[C]{\small\color{gray}\thepage}
\fancyfoot[R]{\small\color{gray}\itshape 大队长出品}
\renewcommand{\headrulewidth}{0.3pt}
\renewcommand{\headrule}{\color{teal3!40}\hrule width\headwidth height\headrulewidth}

% ---- 水印（每页对角线「大队长出品」）----
\AddToShipoutPictureBG{%
  \ifnum\value{page}>1
    \put(\LenToUnit{0.20\paperwidth},\LenToUnit{0.45\paperheight}){%
      \rotatebox{45}{\color{gray!8}\fontsize{120}{120}\selectfont\bfseries 大队长出品}%
    }%
  \fi
}

% ---- 标题样式 ----
\titleformat{\section}
  {\Large\bfseries\color{teal1}}
  {\thesection}{0.6em}{}
  [\vspace{0.2em}{\color{teal3!50}\hrule height 0.5pt}]

\titleformat{\subsection}
  {\large\bfseries\color{teal2}}
  {\thesubsection}{0.5em}{}

\titleformat{\subsubsection}
  {\normalsize\bfseries\color{teal2}}
  {\thesubsubsection}{0.4em}{}

% ---- 段落 ----
\setlength{\parindent}{0pt}
\setlength{\parskip}{0.75em}
\linespread{1.38}

% ---- 列表间距优化 ----
\setlist{itemsep=0.35em, parsep=0.25em, topsep=0.4em}

% ---- 题目级标题 (subsubsubsection 在 pandoc 用 \paragraph) 额外上下间距 ----
\titleformat{\paragraph}
  {\normalsize\bfseries\color{teal1}}
  {\theparagraph}{0.4em}{}
\titlespacing*{\paragraph}
  {0pt}{1.4em plus 0.3em minus 0.2em}{0.6em}

% ---- 节、小节标题上下额外联试分割感 ----
\titlespacing*{\section}
  {0pt}{2.2em plus 0.4em minus 0.2em}{1.0em}
\titlespacing*{\subsection}
  {0pt}{1.6em plus 0.3em minus 0.2em}{0.7em}
\titlespacing*{\subsubsection}
  {0pt}{1.2em plus 0.3em minus 0.2em}{0.5em}

% ---- 数学卷封面 ----
\newcommand{\makemathcoverpage}{%
  \begin{titlepage}
    \thispagestyle{empty}
    \pagecolor{teal1}
    \begin{center}
      \vspace*{2.5cm}
      {\color{white!80}\rule{6cm}{0.6pt}\par}
      \vspace{1cm}
      {\fontsize{52}{60}\selectfont\bfseries\color{white} 算法的低语\par}
      \vspace{0.6cm}
      {\fontsize{30}{36}\selectfont\bfseries\color{teal3!40!white} 数学卷·习题伴读\par}
      \vspace{0.4cm}
      {\fontsize{16}{20}\selectfont\itshape\color{teal3!60!white} 从看懂到会用\par}
      \vspace{1cm}
      {\color{white!80}\rule{6cm}{0.6pt}\par}

      \vspace{2.5cm}
      \begin{minipage}{0.72\textwidth}
        \itshape\color{white!75}\small\setlength{\parskip}{0.5em}
        本册是《算法的低语·数学卷》的姊妹册。主册负责把从香农熵到扩散模型的每个公式讲明白，伴读负责让读者亲手算出来。

        全书 52 节与主册节号逐一对应。每节提供：

        \begin{itemize}\setlength\itemsep{0.2em}\color{white!75}
          \item \textbf{公式卡片}：该节必记公式·变量释义·维度表
          \item \textbf{轻·中·重三档题}：从口算代入到纸笔全推
          \item \textbf{完整推导解答}：每步为什么能走、用了哪条定理
          \item \textbf{错点与陷阱提示}：看不见的踩坑经验
        \end{itemize}

        适合与主册配套阅读，也可独立作为推导示例集使用。
      \end{minipage}

      \vfill
      {\color{white}\Large\bfseries 大队长\enspace 出品\par}
      \vspace{0.3em}
      {\color{teal3!50!white}\small\href{mailto:fqsx@mail.ustc.edu.cn}{fqsx@mail.ustc.edu.cn}\par}
      \vspace{0.6em}
      {\color{teal3!40!white}\small\bfseries v1.2 全册 \,·\, 51 节·汇编完整版\par}
      \vspace{0.5cm}
    \end{center}
    \AddToShipoutPictureFG*{%
      \AtPageUpperLeft{%
        \put(40,-50){\color{white!50}\scriptsize\itshape 大队长出品}%
      }%
      \AtPageUpperLeft{%
        \put(\LenToUnit{\paperwidth-3.2cm},-30){\color{teal3!50!white}\scriptsize\bfseries v1.6}%
      }%
    }
    \clearpage
    \nopagecolor
  \end{titlepage}
}
